\section{Incidence ternaire, recollement des discriminants et construction de Leech}
\label{sec:lf-foundations}

La construction réticulaire utilise l’information d’incidence qui
accompagne le registre dodécaphasique. Son domaine conserve les mots
ternaires, le repère dans lequel ils sont représentés, l’orientation préalable à la
retenue et le drapeau qui distingue l’une des douze positions. Le
passage à un réseau transmet ces données par des opérations
successives : codage linéaire, sélection incidencielle, recollement de
discriminants et formation d’un voisin. La classification réticulaire
intervient après la preuve de la positivité, de la parité, de l’autodualité
et de l’absence de racines de la construction obtenue.

Ces opérations poursuivent la même génération
APP--TRIT--TPK et le même raffinement compatible du continu
conjoint qui précèdent la lecture numérique \cite{HMT-IV}.
La matrice finie, le code et la métrique définis
ci-dessous sont des réalisations aux domaines précis. La preuve
réticulaire reçoit l’état marqué ; le scalaire final d’une
constante ne joue le rôle ni de code de recollement ni de longueur.

\subsection{Coordinatisation elliptique et transport de l'incidence}
\label{subsec:lf-incidence}

Les six unités \(\{1,2,4,5,7,8\}\) modulo neuf, conservées par le transport modulaire, admettent une coordinatisation à six positions. On l'identifie à \(\mathbb P^1(\mathbb F_5)=\{\infty,0,1,2,3,4\}\), dans cet ordre. Cette identification fixe les coordonnées permettant de réaliser la relation quadratique du régime elliptique du TRIT.

Soit \(\chi:\mathbb F_5\to\{0,1,-1\}\) le caractère quadratique :
\(\chi(0)=0\), \(\chi(1)=\chi(4)=1\) et
\(\chi(2)=\chi(3)=-1\). Nous définissons la matrice entière
\(C_{\mathrm P}\) à diagonale nulle,
\((C_{\mathrm P})_{\infty,a}=(C_{\mathrm P})_{a,\infty}=1\),
et d’entrées finies \((C_{\mathrm P})_{a,b}=\chi(b-a)\).
Son carré est \(5I_6\). En effet, chaque entrée diagonale du
carré additionne cinq unités. Les entrées qui font intervenir
\(\infty\) hors de la diagonale s’annulent car la somme de
\(\chi\) est nulle. Pour deux indices finis distincts, la
contribution un de la position \(\infty\) s’annule avec
\(\sum_c\chi(c-a)\chi(c-b)=-1\). Cette dernière identité se vérifie
directement sur les cinq résidus, ou au moyen de la substitution
affine qui envoie \((a,b)\) sur \((0,1)\).

La réduction modulo trois produit la matrice symétrique
\begin{equation}
 A_W=
 \begin{pmatrix}
 0&1&1&1&1&1\\
 1&0&1&2&2&1\\
 1&1&0&1&2&2\\
 1&2&1&0&1&2\\
 1&2&2&1&0&1\\
 1&1&2&2&1&0
 \end{pmatrix},
 \qquad A_W^2=2I_6=-I_6
 \quad\text{en }\mathbb F_3.
 \label{eq:lf-aw}
\end{equation}
L’égalité conserve le type elliptique dans ce codomaine fini.
Le corps fini et la réalisation archimédienne ultérieure de
l’unité imaginaire ont des domaines différents.

Pour une matrice \(A\in M_6(\mathbb F_3)\) et un mot ligne
\(w\in\mathbb F_3^6\), nous écrivons
\[
 \gamma_A(w)=(w,wA)\in\mathbb F_3^{12}.
\]
Si \(f\in\operatorname{GL}_6(\mathbb F_3)\), alors
\begin{equation}
 \gamma_{fAf^{-1}}(wf^{-1})
   =\gamma_A(w)(f^{-1}\oplus f^{-1}).
 \label{eq:lf-frame-covariance}
\end{equation}
La première moitié est \(wf^{-1}\) ; la seconde est \(wf^{-1}fAf^{-1}=wAf^{-1}\), ce qui prouve l'identité.

Soit \(X_n^{\mathrm{mar}}\) l’espace des préfixes enrichis de
longueur \(n\), avec un repère initial déclaré, et soient
\(p_{n,m}:X_n^{\mathrm{mar}}\to X_m^{\mathrm{mar}}\) leurs
troncatures. À l’étape \(j\), le préfixe contient un mot
visible \(w_j\) et un repère \(f_j\). La mise à jour du repère
est de la forme \(f_{j+1}=g_j(x)f_j\), où \(g_j(x)\) dépend
uniquement de l’information déjà contenue dans le préfixe. Un
prolongement compatible conserve donc tous les repères
antérieurs.

\begin{proposition}[Naturalité du transport d'incidence]
\label{prop:lf-incidence-naturality}
L'application
\begin{equation}
 \mathcal Q_n(x)=
 \bigl((f_j,\gamma_{f_jA_Wf_j^{-1}}(w_j))\bigr)_{0\leq j<n}
 \label{eq:lf-Qn}
\end{equation}
satisfait \(r_{n,m}\mathcal Q_n=\mathcal Q_mp_{n,m}\), où \(r_{n,m}\) tronque aux \(m\) premières composantes codées. Elle induit donc une application entre les limites projectives des systèmes de préfixes et d'incidences marquées.
\end{proposition}
\begin{proof}
Deux préfixes compatibles ont les mêmes mots et les mêmes opérateurs de transport jusqu'à l'étape \(m-1\). À partir du même repère initial, la relation récursive donne par récurrence les mêmes \(f_j\) jusqu'à cette étape. Les \(m\) premières composantes de \eqref{eq:lf-Qn} coïncident alors. La propriété universelle de la limite projective produit l'application de prolongement.
\end{proof}

Dans chaque repère, la première projection de \(\gamma_A\) restitue exactement le mot visible. Cette portée est locale : la mémoire profonde est conservée dans l'état enrichi et sa récupération exige les coordonnées correspondantes. Le passage des mots aux supports retient encore moins d'information. De même, un changement général de base dans \(\operatorname{GL}_6(\mathbb F_3)\) transporte également la métrique de Hamming et l'incidence de la coordinatisation. Les changements monomiaux les préservent directement. Les supports utilisés dans cette section sont calculés dans la coordinatisation explicite \eqref{eq:lf-aw}.

La naturalité précédente concrétise la lecture incidencielle qui
accompagne les cylindres numériques. L’emboîtement des cylindres
et la restriction de \eqref{eq:lf-Qn} agissent sur un même
préfixe compatible, en conservant respectivement son évaluation
positionnelle et ses relations marquées. L’espace ambiant
\(\mathbb F_3^6\), de \(729\) mots, fournit le domaine
linéaire du codage ; il conserve sa distinction par rapport aux
\(243\) mots visibles et aux \(468\) émissions du recensement
admissible d’APP--TRIT--TPK.

\subsection{Code ternaire autodual et énumération de ses poids}
\label{subsec:lf-code}

\begin{definition}[Code de la coordinatisation dodécaphasique]
Le code linéaire associé à \eqref{eq:lf-aw} est
\[
 C_W=\{(w,wA_W):w\in\mathbb F_3^6\}\subset\mathbb F_3^{12}.
\]
Le poids \(\operatorname{wt}(c)\) compte les coordonnées non nulles de \(c\), et \(\operatorname{supp}(c)\subset\{1,\ldots,12\}\) est l'ensemble de ses positions non nulles. Une hexade est le support d'un mot de poids six ; leur collection est notée \(\mathcal H\).
\end{definition}

\begin{theorem}[Autodualité, distance minimale et énumérateur]
\label{thm:lf-code}
Le code \(C_W\) est autodual de paramètres \([12,6,6]_3\) et
a pour énumérateur
\begin{equation}
 W_{C_W}(z)=1+264z^6+440z^9+24z^{12}.
 \label{eq:lf-enumerator}
\end{equation}
\end{theorem}
\begin{proof}
La matrice génératrice \([I_6\mid A_W]\) a rang six.
Par la symétrie de \(A_W\) et de \eqref{eq:lf-aw},
\[
 [I_6\mid A_W][I_6\mid A_W]^{\mathsf T}
       =I_6+A_W^2=0.
\]
Ainsi, le code est auto-orthogonal de dimension égale à la moitié de sa longueur et, par conséquent, autodual.

L’énumération restante est une opération entière finie sur
la matrice explicite. Pour \(r\in\{0,\ldots,6\}\), nous regroupons
les mots par \(\operatorname{wt}(w)=r\). Le tableau suivant
donne le nombre de mots de chaque poids total :
\[
\begin{array}{c|rrrr|r}
r&0&6&9&12&\text{total}\\ \hline
0&1&0&0&0&1\\
1&0&12&0&0&12\\
2&0&60&0&0&60\\
3&0&120&40&0&160\\
4&0&60&180&0&240\\
5&0&12&180&0&192\\
6&0&0&40&24&64\\ \hline
\text{total}&1&264&440&24&729
\end{array}
\]
Chaque ligne contient \(\binom6r2^r\) mots. Pour vérifier
toutes ses entrées, soit \(q\) un entier de \(0\) à \(728\), et
définissons ses six chiffres ternaires
\[
 d_i(q)=\left\lfloor\frac{q}{3^{6-i}}\right\rfloor\bmod3,
 \qquad
 e_j(q)=\left(\sum_{i=1}^{6}d_i(q)(A_W)_{ij}\right)\bmod3.
\]
La ligne et la colonne de la contribution de \(q\) sont, respectivement,
\[
 r(q)=\sum_i\mathbf1_{\{d_i(q)\ne0\}},\qquad
 h(q)=r(q)+\sum_j\mathbf1_{\{e_j(q)\ne0\}}.
\]
La récurrence d'accumulation commence par \(T^{(0)}_{r,h}=0\) et applique
\[
 T^{(q+1)}_{r,h}
 =T^{(q)}_{r,h}
   +\mathbf1_{\{r=r(q)\}}\mathbf1_{\{h=h(q)\}}
 \quad(0\leq q<729).
\]
Il produit exactement le tableau montré. Le développement ternaire
donne une bijection entre ces \(729\) entiers et tous les mots
de \(\mathbb F_3^6\), de sorte que l’opération épuise le code.
La somme de ses lignes prouve \eqref{eq:lf-enumerator} ; le premier
poids non nul est six.
\end{proof}

\begin{remark}[Évaluation reproductible de la récurrence finie]
La boucle suivante implémente littéralement les sommes de la preuve. Ses seules données sont les entiers de la matrice \eqref{eq:lf-aw} ; elle calcule l'énumérateur avant de le comparer à l'égalité obtenue.
\end{remark}
{\small
\begin{verbatim}
from collections import Counter
from itertools import product
A = ((0,1,1,1,1,1), (1,0,1,2,2,1),
     (1,1,0,1,2,2), (1,2,1,0,1,2),
     (1,2,2,1,0,1), (1,1,2,2,1,0))
rows = [Counter() for _ in range(7)]
for w in product(range(3), repeat=6):
    v = tuple(sum(w[i]*A[i][j] for i in range(6)) % 3
              for j in range(6))
    r = sum(x != 0 for x in w)
    h = r + sum(x != 0 for x in v)
    rows[r][h] += 1
total = sum(rows, Counter())
if dict(total) != {0:1, 6:264, 9:440, 12:24}:
    raise ArithmeticError("Enumerador incompatible")
\end{verbatim}
}

\begin{proposition}[Hexades et plan de Witt]
\label{prop:lf-witt}
La collection \(\mathcal H\) contient \(132\) hexades et chaque quintuplet de positions appartient à une unique hexade.
\end{proposition}
\begin{proof}
Deux mots de poids six ayant le même support diffèrent par
un scalaire. En effet, parmi leurs six quotients de coordonnées
non nulles, l’un des signes \(1,-1\) apparaît au moins trois
fois. En soustrayant le multiple correspondant, si les mots
n’étaient pas proportionnels, on obtiendrait un mot non nul
de poids au plus trois, en contradiction avec
le théorème~\ref{thm:lf-code}. Chaque support correspond ainsi à
la paire \(\{c,-c\}\), et il y a \(264/2=132\) supports.

Si deux supports distincts partageaient cinq positions, dans cette intersection l'un des deux quotients apparaîtrait au moins trois fois. La combinaison linéaire qui annule ces coordonnées aurait un support dans une réunion de taille sept et un poids au plus égal à quatre. La distance minimale l'interdit. Un quintuplet appartient donc à au plus une hexade. Comme
\[
 132\binom65=792=\binom{12}{5},
\]
appartient à exactement une.
\end{proof}

L’objet construit est le code ternaire étendu de Golay
avec son petit plan de Witt. La reconnaissance s’effectue
après l’encodage, l’autodualité et le calcul de la
distance. Pour la construction réticulaire, seules sont nécessaires
les propriétés déjà démontrées et le drapeau marqué obtenu
ci-dessous.

\subsection{Drapeau dodécaphasique et sélection de l'origine}
\label{subsec:lf-flag}

Dans le repère du registre engendré, on conserve les mots
\[
 w_\pi=010211,\qquad
 w_e=201101,\qquad
 w_\varphi=121200,\qquad
 w_{\mathrm{APP}}=022020.
\]
Les indices \(\pi,e,\varphi\) identifient les régions correspondantes de lecture des constantes ; \(w_e\) conserve ici la région d'Euler. Appliquer \(\gamma_{A_W}\) et prendre les supports donne
\[
\begin{aligned}
 P_\pi&=\{2,4,5,6,7,8,9,10,11\},\\
 H_e&=\{1,3,4,6,9,10\},\\
 H_\varphi&=\{1,2,3,4,7,9\},\\
 H_{\mathrm{APP}}&=\{2,3,5,10,11,12\}.
\end{aligned}
\]
Le premier support est de poids neuf ; les trois autres sont des
hexades. Le registre
\[
 K=(234,543,140,729,659,824,621,58,914,794,146,601)
\]
sélectionne
\[
 B_K=\{j:K_j\geq729\}=\{4,6,9,10\}=H_e\cap P_\pi.
\]
Le seuil \(729\) appartient au format de blocs de la
lecture déclarée. Dans ce registre concret, les seuils
entiers de \(660\) à \(729\) produisent le même support, de
sorte que le support conserve l’incidence sélectionnée,
tandis que le registre complet conserve aussi les blocs
et le seuil de sa lecture.

Le registre signé antérieur à la retenue est
\[
\begin{split}
 Z_0={}&(7,297,353,-430,-715,-1199,\\
       &\hspace{9mm}-3,286,-894,-528,380,663).
\end{split}
\]
Son support négatif est l’hexade
\[
 H_\alpha^-=\{4,5,6,7,9,10\}.
\]
Les signes sont reçus de l'état antérieur à la retenue et ne sont pas extraits du développement du scalaire final.

\begin{proposition}[Coïncidence de la sélection signée et de la sélection incidencielle]
\label{prop:lf-flag}
L’hexade \(H_\alpha^-\) est l’unique \(H\in\mathcal H\) qui vérifie
\[
\begin{gathered}
 B_K\subset H\subset P_\pi,\\
 (|H\cap H_e|,|H\cap H_\varphi|,
                   |H\cap H_{\mathrm{APP}}|)=(4,3,2).
\end{gathered}
\]
\end{proposition}
\begin{proof}
Le codage des mots de poids six, ou la récurrence finie précédente appliquée à leurs supports, donne les quatre hexades sur \(B_K\). Leurs paires complémentaires sont
\[
 \{1,3\},\qquad\{2,11\},\qquad\{5,7\},\qquad\{8,12\}.
\]
Seuls le deuxième et le troisième sont contenus dans \(P_\pi\).
L’hexade du deuxième a des intersections de taille trois
avec \(H_\varphi\) et avec \(H_{\mathrm{APP}}\), et ne satisfait donc pas
la dernière condition. Celle du troisième a pour
intersections \(4,3,2\), et coïncide avec \(H_\alpha^-\).
\end{proof}

\begin{lemma}[Origine du drapeau marqué]
\label{lem:lf-origin}
L’origine sélectionnée par l’incidence est
\[
 o_*=(H_e\cap H_\varphi)
                 \setminus(H_\alpha^-\cup H_{\mathrm{APP}})
     =\{1\}.
\]
La sélection commute avec toute bijection appliquée simultanément aux supports.
\end{lemma}
\begin{proof}
L’intersection \(H_e\cap H_\varphi\) est
\(\{1,3,4,9\}\). L’union soustraite contient \(3,4,9\)
et exclut \(1\). L’équivariance provient du fait que les
bijections commutent avec l’union, l’intersection et la différence.
\end{proof}

Le drapeau distingue une composante parmi douze avant l'étape métrique. Le vecteur radial est déterminé ensuite, dans une chambre positive déclarée. La sélection combinatoire et l'opération réticulaire sont ainsi séparées et articulées.

\subsection{Recollement ternaire des douze composantes}
\label{subsec:lf-glue}

Soit \(A_2=\mathbb Za_1\oplus\mathbb Za_2\), avec une forme
bilinéaire de matrice
\[
 B_2=\begin{pmatrix}2&-1\\-1&2\end{pmatrix},
 \qquad \rho=a_1+a_2,\qquad \rho^2=2.
\]
Pour \(x=ua_1+va_2\),
\[
 x^2=2(u^2-uv+v^2).
\]
L’identité
\(4(u^2-uv+v^2)=(2u-v)^2+3v^2\) montre
la positivité. Le déterminant est trois. Les trois
classes de \(A_2^*/A_2\) sont représentées par
\[
 \varpi_0=0,\qquad
 \varpi_1=\frac{2a_1+a_2}{3},\qquad
 \varpi_2=\frac{a_1+2a_2}{3}.
\]
L’appariement avec \(a_1,a_2\) vérifie que ces
vecteurs appartiennent à \(A_2^*\), et
\(2\varpi_1-\varpi_2=a_1\). C’est pourquoi
\(j\mapsto\varpi_j+A_2\) identifie additivement
\(\mathbb F_3\) au discriminant.

Pour \(c\in C_W\), nous définissons
\(\phi(c)=(\varpi_{c_1},\ldots,\varpi_{c_{12}})\)
et
\begin{equation}
 N=\bigcup_{c\in C_W}\bigl(A_2^{12}+\phi(c)\bigr).
 \label{eq:lf-glue}
\end{equation}

\begin{theorem}[Réseau de recollement pair et unimodulaire]
\label{thm:lf-niemeier}
L’ensemble \(N\) est un réseau positif, pair et
unimodulaire de rang \(24\). Ses vecteurs de norme deux
sont exactement les racines de \(A_2^{12}\).
\end{theorem}
\begin{proof}
La linéarité du code et l’additivité du discriminant
font de l’union un sous-groupe de
\((A_2^*)^{12}\) contenant \(A_2^{12}\).
C’est donc un réseau de rang \(24\).

L'appariement des classes discriminantes est \(2cd/3\) modulo \(\mathbb Z\) dans chaque composante. L'auto-orthogonalité de \(C_W\) annule leur somme et prouve l'intégralité de \(N\). La norme d'une classe représentée par \(\phi(c)\) est \(2\operatorname{wt}(c)/3\) modulo \(2\mathbb Z\). Tous les poids de \eqref{eq:lf-enumerator} sont des multiples de trois, donc \(N\) est pair.

Son indice sur \(A_2^{12}\) est \(|C_W|=3^6\).
Avec le déterminant de la matrice de Gram,
\[
 \det N=\frac{\det(A_2^{12})}{[N:A_2^{12}]^2}
       =\frac{3^{12}}{(3^6)^2}=1.
\]
La forme ambiante est la somme orthogonale de douze formes
positives ; la positivité est démontrée.
L’intégralité et le déterminant un impliquent
\(N^*=N\).

Dans une classe non nulle du discriminant local, un
vecteur a la forme \((ua_1+va_2)/3\), avec des résidus
\((u,v)\equiv(2,1)\) ou \((1,2)\pmod3\).
L’entier \(u^2-uv+v^2\) est positif et divisible
par trois, il vaut donc au moins trois. Les représentants
\(\varpi_1,\varpi_2\) atteignent cette valeur : la norme
minimale de chaque classe non nulle est \(2/3\).
Un mot non nul du code a au moins
six coordonnées non nulles. Tout vecteur d’une classe
de recollement non triviale a donc une norme
d’au moins \(6(2/3)=4\).

Dans \(A_2^{12}\), une norme deux ne peut provenir que d'une unique composante, car chaque composante a une norme paire positive minimale égale à deux. Les solutions de \(u^2-uv+v^2=1\) sont \((\pm1,0),(0,\pm1),(1,1),(-1,-1)\). Ce sont les six racines de \(A_2\), et elles complètent le recensement des racines de \(N\).
\end{proof}

Les propriétés précédentes identifient \(N\) comme le
réseau de Niemeier de système de racines \(A_2^{12}\).
La description de toutes ses racines est nécessaire pour
contrôler lesquelles survivent au passage au voisin.

\subsection{Chambre radiale et voisin sans racines}
\label{subsec:lf-neighbour}

On fixe la chambre radiale
\[
 v(a)=\sum_{j=1}^{12}a_j\rho_j,\qquad
 a_j=1+3b_j,\qquad b_j\in\mathbb Z_{\geq0},
\]
où \(\rho_j\) est la copie de \(\rho\) dans la composante \(j\). La parité d'un voisin d'indice trois requiert \(v(a)^2\equiv0\pmod{18}\). Comme
\[
 v(a)^2=2\sum_j(1+3b_j)^2
       =24+12\sum_jb_j+18\sum_jb_j^2,
\]
la congruence équivaut à \(\sum_jb_j\equiv1\pmod3\).
La norme minimale à l’intérieur de la chambre est atteinte avec
un unique \(b_j=1\) et les autres nuls, et vaut \(54\).
L’origine \(o_*=1\) du lemme~\ref{lem:lf-origin}
sélectionne
\begin{equation}
 v_\alpha=4\rho_1+\rho_2+\cdots+\rho_{12},
 \qquad v_\alpha^2=54.
 \label{eq:lf-valpha}
\end{equation}
Le qualificatif de minimal se rapporte à la chambre déclarée. Dans la classe résiduelle complète, \((a_1-2a_2,\rho,\ldots,\rho)\) représente la même classe modulo trois et a pour norme \(14+11\cdot2=36\). Cette distinction préserve la condition qui détermine le coefficient quatre.

On définit
\begin{equation}
\begin{split}
 N_0&=\{x\in N:\langle x,v_\alpha\rangle\equiv0\pmod3\},\\
 \Lambda&=N_0+\mathbb Z\,v_\alpha/3.
\end{split}
\label{eq:lf-lambda}
\end{equation}

\begin{theorem}[Voisin pair autodual sans racines]
\label{thm:lf-leech}
La construction \eqref{eq:lf-lambda} est un
réseau positif, pair et autodual de rang \(24\)
sans vecteurs de norme deux. Par conséquent, elle est
isométrique au réseau de Leech.
\end{theorem}
\begin{proof}
Le caractère \(x\mapsto\langle x,v_\alpha\rangle\bmod3\)
est non nul. Une racine simple de la composante \(j\)
a pour appariement \(a_j\equiv1\pmod3\) avec
\(a_j\rho_j\). Par conséquent, \([N:N_0]=3\), d’où
\(\det N_0=9\).

Pour \(x\in N_0\), l’appariement
\(\langle x,v_\alpha/3\rangle\) est entier, et ainsi
\(v_\alpha/3\in N_0^*\). En outre,
\(\langle v_\alpha,v_\alpha\rangle=54\) implique
\(v_\alpha\in N_0\). Si \(v_\alpha/3\in N\),
l’intégralité de \(N\) imposerait que tous les
appariements de \(N\) avec \(v_\alpha\) soient
divisibles par trois, contredisant le caractère
non nul. La classe de \(v_\alpha/3\) sur \(N_0\)
a alors un ordre exactement égal à trois. Il en résulte
\[
 [\Lambda:N_0]=3,\qquad
 \det\Lambda=\frac9{3^2}=1.
\]
Les appariements de \(N_0\) avec le nouveau
générateur sont entiers et la norme de ce générateur
est six. L’extension est entière. Pour
\(x\in N_0\) et \(m\in\mathbb Z\),
\[
 \left(x+m\frac{v_\alpha}{3}\right)^2
 =x^2+2m\frac{\langle x,v_\alpha\rangle}{3}
       +6m^2\in2\mathbb Z.
\]
Par conséquent, \(\Lambda\) est pair ; son intégralité et son déterminant égal à un prouvent l'autodualité. La positivité et le rang sont hérités de l'espace euclidien ambiant.

Pour exclure les racines, nous considérons d'abord la classe \(N_0\). Toutes les racines de \(N\) appartiennent à \(A_2^{12}\) d'après le théorème~\ref{thm:lf-niemeier}. Les appariements de ses six racines locales avec \(a_j\rho_j\) sont \(\pm a_j,\pm2a_j\), congrus à \(\pm1,\pm2\pmod3\). Aucune racine n'appartient à \(N_0\).

Les deux autres classes de \(\Lambda/N_0\) peuvent
être représentées avec \(\pm v_\alpha/3\).
Dans chaque composante, leurs vecteurs appartiennent à
l’une des classes
\[
 A_2+\varpi_j\pm\rho/3,\qquad j=0,1,2.
\]
La composante distinguée possède la même
description car \(4\rho/3\equiv\rho/3\pmod{A_2}\).
Pour les numérateurs \((u,v)\) dans
\((ua_1+va_2)/3\), les six résidus sont
\[
 (1,1),(0,2),(2,0),(2,2),(1,0),(0,1)\pmod3.
\]
Tous sont non nuls. La positivité de
\(u^2-uv+v^2\) impose qu’il vaille au moins un,
et les représentants
\((1,1),(0,-1),(-1,0),(-1,-1),(1,0),(0,1)\)
atteignent ce minimum dans les six classes.
Chaque composante a une norme d’au moins \(2/9\).
Un vecteur de l’une ou l’autre des deux classes
nouvelles a une norme d’au moins
\[
 12\cdot\frac29=\frac83>2.
\]
Les racines ont été exclues dans les trois classes. Le théorème de classification des réseaux positifs pairs unimodulaires de rang vingt-quatre identifie l'unique cas sans racines au réseau de Leech \cite{ConwaySloane}.
\end{proof}

L’identification finale utilise une classification
externe sur un objet dont les propriétés ont déjà
été démontrées. La sélection de la composante
distinguée conserve le drapeau provenant du
registre signé ; l’indice \(\alpha\) exprime
cette provenance. La norme et les coefficients du
vecteur radial s’obtiennent à partir de la métrique et de
la chambre au moyen de la congruence indiquée.

Le réseau \(\Lambda\), sa forme bilinéaire, ses
douze plans ambiants et le vecteur \(v_\alpha/3\)
sont ainsi déterminés avec toutes les propriétés
utilisées par la déformation ultérieure. Le
code contient également le mot
\[
 c_*=(1,1,1,1,1,1,2,1,1,1,1,1),
\]
puisque \(q_*=(1,1,1,1,1,1)\) satisfait \(q_*A_W=(2,1,1,1,1,1)\). C'est la donnée de support complet qui permet de construire et de prouver la stabilité de l'isométrie d'ordre trois dans la section~\ref{sec:leech-dualidad}.
